\subsection{Componentes funcionales y separación de responsabilidades}
Las cuentas 1.4, 1.5 y 1.6 contienen 30 componentes funcionales y 306 paquetes terminales. La cuenta 1.4 desarrolla Operación marítima, Escuela náutica e Infraestructura; 1.5, Varadero y contratistas, Combustible y Ambiente y consumos; 1.6, Clientes y contratos y Finanzas y cobranza. Auditoría y seguridad permanece en 1.3. Esta organización conserva los nueve dominios de SD3/SD4/SD5: un componente EDT representa trabajo para entregar una capacidad y no crea otro dominio o autoridad de datos.

Cada componente conserva sus tres paquetes originales y añade lotes estimables de análisis/diseño, construcción, verificación y coordinación. Los 90 códigos originales y sus 4.200 HH se acreditan una vez; la selección UCP suma 12.708 HH en 306 paquetes, con 7.766 HH en E1 y 4.942 HH en E2. Los componentes usan las bases comunes 1.2/1.3/1.7 y los datos migrados 1.8. El diccionario de conciliación identifica las HH adicionales y la matriz de actividades muestra el crédito del trabajo original, sin reconstruir servicios comunes.

Los resultados se vinculan con los ámbitos RF de SD3 y las entidades/autoridades de SD5. La comprobación exhaustiva de cobertura deberá relacionar cada ID aplicable de T-12 con el paquete y su prueba, incluidos requisitos transversales y complementos del caso. Una referencia al ámbito NAV, ESC o FAC no equivale a demostrar por sí sola todos sus requisitos; la matriz conservará IDs y enunciados originales \parencite[caps.~19--20]{bases-transversales}.

E1 mantiene movimientos, planes, alertas, amarra compatible, escuela, inspección de boyas, captura de combustible, medición e historia ambiental inicial, junto con los maestros y contratos que requieren esos procesos. E2 completa agenda y faenas, contratistas, residuos, imputación comercial, espera, cobranza y expedientes. La medición temprana conserva su dependencia de instalación y serie anterior a 2029. Las innovaciones conservan su alcance incremental y sus pilotos de SD13, sin sustituir capacidades obligatorias.
